%% ============================================================
%% LIMITATIONS -- v2 merge, drop-in replacement for <<KEEP:limitations-other>>.
%% Placed AFTER the two paragraphs already drafted in the manuscript
%% (the scope-ceiling paragraph and the pre-upgrade-validation paragraph),
%% which are not repeated here.
%%
%% Source: v1 lines 2085-2157, remapped (v1 L3 -> \Lone, v1 L3+ -> \Lthree,
%% v1 L3NL -> \NLref) and filtered: several v1 limitations have become v2 RESULTS
%% and must NOT be restated as open.
%%
%% DROPPED from v1, because the revision resolved them:
%%  - "L+ -> Lnl interpretation": v1 conceded linearisation and delay could not be
%%    separated. They now are -- delay is its own level (L6, measured at 0.00 %) and
%%    the linearisation error is solved directly. Restating this as a limitation would
%%    contradict Section 4.7 and concede R1.3/R2.3 unnecessarily.
%%  - "Supply temperature ... could amplify the L3->L+ effect": now bounded forward
%%    by the supply-temperature study (Section 4.6), so it is reported as a bounded
%%    result, not an open risk.
%%
%% CLOSES: R1.7 (deterministic, hourly, no reserves) -- previously answered in the
%% response letter with "limitations expanded" while no expansion existed.
%% COMPLETES: R1.5 (precomputed COP and sub-hourly dynamics; the heating-curve half
%% is already covered by the scope-ceiling paragraph above).
%% ============================================================

\paragraph{Uncertainty, reserves and sub-hourly operation.}
The dispatch is deterministic and hourly: prices and demand are perfect-foresight series,
and forecast error, reserve provision, intra-hour excursions and unit outages are excluded.
Each plausibly raises the value of the physics found inert here---transport delay is a
zero-timestep shift at hourly but not five-minute resolution, pipe thermal inertia buffers
precisely the sub-hourly excursions we omit, and storage and heat-pump flexibility are worth
more in reserve markets than against a known price curve. The nulls should thus be read as
decision-irrelevance \emph{in the deterministic hourly planning regime} where most published
dispatch models operate, not as a claim about real-time control; per
Section~\ref{subsec:whynull}, relaxing these assumptions can only give the physics more room,
so the nulls are conservative while the loss-dominance result, independent of them, is not.

A caution runs the other way too. \citet{Koek2026} show that in multistage investment
planning, ignoring price and demand-trajectory uncertainty biases technology selection by
27--39\,\%. Their setting is investment, not dispatch, and fixes topology, but it suggests
the resolution bias quantified here and their uncertainty bias are additive rather than
substitutable: a model can be wrong about the network and about the future independently,
and addressing both in one formulation remains open.

\paragraph{Demand-side flexibility.}
All demands are exogenous time series and demand response is excluded. This is one of the
few exclusions that could plausibly \emph{increase} the topology effect rather than leave
it unchanged, because spatially distributed flexibility has location-dependent value: a
flexible load near the network far end can absorb local surplus without incurring the
trunk losses that serving it from the plant would. Whether that is enough to lift routing
above the decision threshold is untested here and is a well-posed question for a follow-up
study.

\paragraph{Network topology class and generation siting.}
Every case, real and synthetic, is a radial tree with unidirectional flow. Ring and meshed
topologies are not merely a larger instance of the same problem: with pipe-flow direction
fixed by the pipe list, the return-side pressure balance around a cycle forces the sum of
the loop's friction drops to be non-positive, which is infeasible for any non-zero flow.
Admitting loops requires signed-flow hydraulic variables, a per-pipe direction selection
that reorients each pressure drop and temperature mix, which is a distinct formulation
rather than a parameter change. The same boundary limits the generation-topology
moderator: our synthetic generator roots all injection at the primary producer, so the
distributed-generation arm is posed as a controlled open question rather than answered,
and multi-source operation on a meshed network is left to the companion study. Because
generation is central in every case examined, the null we report for spatial routing is
established for centrally-supplied networks specifically.

\paragraph{The nonlinear reference.}
The nonlinear reference is solved on representative weeks rather than globally. The
full-year non-convex programme finds no incumbent within the solver budget, and a low-load
summer week is infeasible under native physics with the linearised schedule fixed. We
report both as results in Section~\ref{subsec:linearisation} (the second is itself an
instance of the physical-deliverability failure the regret evaluation was built to expose),
but they do bound the claim: the solved linearisation error is established for winter
and autumn operating points, and whether it remains sub-percent under shoulder-season
operation with frequent on/off cycling is open. Warm-start heuristics or a reformulation
of the bilinear terms are the natural routes.

\paragraph{One real network, and the synthetic generator's assumptions.}
The Memmingen figures come from a single industrial network; the synthetic factorial
supplies external validity for the \emph{structure} of the result (loss dominance,
adder non-transfer, the loss-number rule) but not for the absolute magnitudes, which
should not be treated as universal constants. The synthetic networks also share the
assumptions of the generator that produced them: radial trees, a common heating curve,
proportionally sized capacity under a stated convention. They are a controlled parametric
family, not a sample of real networks, and the balanced factorial buys unambiguous
attribution across the design at the cost of that representativeness.

\paragraph{Demand-concentration metric.}
The normalised variance index used as a synthetic design factor has a narrow dynamic range
on realistic tree networks, and the model-selection observations in
Table~\ref{tab:criteria} therefore use the Gini coefficient and top-$k$ demand share
instead. An empirical mapping from either metric to the aggregation gap would require
several real networks and is a natural follow-up.
